\section{20. Managing time and pacing}

\subsection{Session length}

For a conventional one-shot, plan approximately four to six hours. Duration depends on encounters fought, familiarity with D\&D and exploration decisions.

\begin{longtable}{p{0.34\linewidth}p{0.60\linewidth}}
\rowcolor{tablehead} \textbf{Phase} & \textbf{Suggested duration} \\
\endfirsthead
\rowcolor{tablehead} \textbf{Phase} & \textbf{Suggested duration} \\
\endhead
Introduction and initial exploration & 20–40 minutes \\
Main exploration & 90–150 minutes \\
Major challenges and treasure discovery & 45–90 minutes \\
Doomsday and escape & 30–60 minutes \\
Conclusion & 10–15 minutes \\
\end{longtable}

These are pacing targets, not mandatory phases. Treasure may be discovered unexpectedly early.

\subsection{Avoiding excessive combat}

Potential encounters can consume an entire session if all are resolved to completion. Encourage retreat through consistent threshold restrictions. Defeating an encounter is not always necessary. Ordinary enemies cannot pursue beyond their chamber, and dangerous rooms without rewards offer little incentive for prolonged combat.

\subsection{Resting}

Established rules impose no specific resting restriction. A recommended survival one-shot default is to allow short rests in suitable safe chambers while discouraging long rests through fictional time pressure. If long rests are prohibited or restricted by magic, establish this before it becomes relevant. Avoid unexpectedly denying recovery abilities players reasonably expected.

\subsection{Maintaining pressure}

Shifting and discovery already provide substantial pressure. No additional countdown is needed unless circumstances warrant it. If players spend excessive time discussing routine choices, summarize options and ask for a decision. Do not arbitrarily punish careful planning; strategic discussion is part of the adventure.
